% META: {"id": "1NSI-TC-CO-013", "chapitre": "1NSI-TYPES-CONSTRUITS", "type_objet": "corrige", "capacites": ["1NSI-TYPES-CONSTRUITS-C1"], "mode_creation": "ex_nihilo", "sources_inspiration": [], "status": "needs_review"}
% BEGIN-VERIFY
% def top_joueurs(joueurs, n):
%     classement = list(joueurs)
%     for i in range(len(classement)):
%         idx_max = i
%         for j in range(i + 1, len(classement)):
%             if classement[j][1] > classement[idx_max][1]:
%                 idx_max = j
%         classement[i], classement[idx_max] = classement[idx_max], classement[i]
%     resultat = []
%     for k in range(min(n, len(classement))):
%         resultat.append(classement[k][0])
%     return tuple(resultat)
% joueurs = [("Alice", 2400), ("Bob", 1800), ("Clara", 2600), ("Dan", 2100)]
% assert top_joueurs(joueurs, 2) == ("Clara", "Alice")
% assert top_joueurs(joueurs, 1) == ("Clara",)
% assert top_joueurs(joueurs, 4) == ("Clara", "Alice", "Dan", "Bob")
% END-VERIFY
\begin{corrige}{1NSI-TC-EX-013}
On utilise un tri par sélection en ordre décroissant sur le score Elo, puis on extrait les \lstinline{n} premiers noms :

\begin{python}
def top_joueurs(joueurs, n):
    # Copie pour ne pas modifier la liste originale
    classement = list(joueurs)
    # Tri par selection (Elo decroissant)
    for i in range(len(classement)):
        idx_max = i
        for j in range(i + 1, len(classement)):
            if classement[j][1] > classement[idx_max][1]:
                idx_max = j
        classement[i], classement[idx_max] = (
            classement[idx_max], classement[i]
        )
    # Extraction des n premiers noms
    resultat = []
    for k in range(min(n, len(classement))):
        resultat.append(classement[k][0])
    return tuple(resultat)
\end{python}

\textbf{Raisonnement :} on copie la liste d'entrée avec \lstinline{list(joueurs)} pour ne pas modifier l'original. Le tri par sélection cherche à chaque étape le tuple ayant le plus grand Elo parmi les éléments restants (\lstinline{classement[j][1]} accède au score Elo du $j$-ème tuple) et le place en position \lstinline{i}. Enfin, on construit le tuple résultat en ne conservant que les noms (\lstinline{classement[k][0]}).

\textbf{Complexité :} le tri par sélection est en $O(n^2)$ où $n$ est le nombre de joueurs. L'extraction finale est en $O(n)$.

Vérification : \lstinline{top_joueurs(joueurs, 2)} renvoie bien \lstinline{("Clara", "Alice")}.
\end{corrige}
